\documentclass[11pt,a4paper]{article}
\usepackage{amsmath,amssymb}
\usepackage[margin=2.5cm]{geometry}
\usepackage{enumitem}
\usepackage{hyperref}
\usepackage{fancyhdr}
\pagestyle{fancy}
\fancyhf{}
\rhead{ETH Zurich Executive Education}
\lhead{Section 6 --- Agentic AI}
\cfoot{\thepage}

\title{Calibration Quiz: Agentic AI\\[0.2em]
       \large What Works, What Breaks, What You Can't Trust}
\author{Companion to the ETH Zurich Executive Lecture}
\date{}

\begin{document}
\maketitle

\noindent\textbf{Instructions.}
Answer each question with a single letter, number, or short phrase.
For forecasting questions, also write a 90\% confidence interval
(a lower and upper bound you are 90\% sure contains the true answer).
This is a calibration exercise --- the aim is not to be right, but
to learn how wide your uncertainty really is.

\vspace{1em}

\section*{Part I --- Facts}

\begin{enumerate}[label=Q\arabic*.]

\item The single most replicated finding in the agentic-AI literature is
      that the main driver of agent performance is:
\begin{enumerate}[label=(\alph*)]
  \item the base model's parameter count
  \item the scaffolding (planner, tools, reflection) around the model
  \item the training data used during pre-training
  \item the choice of prompt template
\end{enumerate}

\item Yao et al.'s \emph{ReAct} paper was published at:
\begin{enumerate}[label=(\alph*)]
  \item NeurIPS 2022 \qquad (b) ICLR 2023 \qquad (c) ICML 2023 \qquad (d) ACL 2023
\end{enumerate}

\item On the ALFWorld household-task benchmark, ReAct raised the same
      base model's success rate from roughly:
\begin{enumerate}[label=(\alph*)]
  \item 10\% to 30\% \qquad (b) 34\% to 71\% \qquad (c) 55\% to 82\% \qquad (d) 60\% to 95\%
\end{enumerate}

\item Schick et al.'s \emph{Toolformer} paper makes which headline claim?
\begin{enumerate}[label=(\alph*)]
  \item A 13B model matches GPT-4 on tool-use benchmarks.
  \item A 6.7B model that teaches itself to call APIs outperforms a
        175B GPT-3 model without tools on downstream tasks.
  \item Tool use is unnecessary once the base model crosses 70B
        parameters.
  \item Language models cannot reliably learn to call external tools
        without reinforcement learning from human feedback.
\end{enumerate}

\item The four ICLR 2024 benchmarks covered in this lecture are GAIA,
      SWE-bench, WebArena, and:
\begin{enumerate}[label=(\alph*)]
  \item BIG-bench \qquad (b) HELM \qquad (c) AgentBench \qquad (d) MMLU
\end{enumerate}

\item On GAIA, the approximate human baseline is:
\begin{enumerate}[label=(\alph*)]
  \item 72\% \qquad (b) 82\% \qquad (c) 92\% \qquad (d) 98\%
\end{enumerate}

\item OpenAI stopped reporting SWE-bench Verified scores for its
      frontier models because:
\begin{enumerate}[label=(\alph*)]
  \item Scores plateaued below 50\%.
  \item Scores became too expensive to compute.
  \item All tested frontier models showed training-data contamination
        on the benchmark.
  \item The benchmark was deprecated by its authors.
\end{enumerate}

\item ``Faith and Fate'' (Dziri et al., NeurIPS 2023) reports that
      GPT-4's correct rate on 5-digit multiplication is approximately:
\begin{enumerate}[label=(\alph*)]
  \item 4\% \qquad (b) 23\% \qquad (c) 59\% \qquad (d) 88\%
\end{enumerate}

\item Langosco et al.'s \emph{Goal Misgeneralization in Deep
      Reinforcement Learning} was published at:
\begin{enumerate}[label=(\alph*)]
  \item NeurIPS 2021 \qquad (b) ICLR 2022 \qquad (c) ICML 2022 \qquad (d) ICML 2023
\end{enumerate}

\item Greshake et al.'s \emph{Not What You've Signed Up For} paper on
      indirect prompt injection was published as:
\begin{enumerate}[label=(\alph*)]
  \item Main-track paper at CCS 2023
  \item Workshop paper at ACM AISec 2023 (co-located with CCS)
  \item Main-track paper at USENIX Security 2023
  \item Journal article in IEEE Security \& Privacy 2023
\end{enumerate}

\item Park et al.'s \emph{AI Deception} survey appeared in which
      peer-reviewed journal?
\begin{enumerate}[label=(\alph*)]
  \item Nature \qquad (b) Science \qquad (c) Patterns (Cell Press) \qquad (d) PNAS
\end{enumerate}

\item Which of the following is \emph{not} one of the four
      peer-reviewed failure modes discussed in this lecture?
\begin{enumerate}[label=(\alph*)]
  \item Compounding errors on long-horizon tasks
  \item Goal misgeneralization
  \item Indirect prompt injection
  \item Hallucinated citations in scientific writing
\end{enumerate}

\item Which of the following is explicitly \emph{not} peer-reviewed
      (a preprint or technical report)?
\begin{enumerate}[label=(\alph*)]
  \item Reflexion (Shinn et al.)
  \item GAIA (Mialon et al.)
  \item Sleeper Agents (Hubinger et al., Anthropic)
  \item Generative Agents (Park et al.)
\end{enumerate}

\end{enumerate}

\section*{Part II --- Forecasting (with 90\% confidence intervals)}

\noindent For each question, give your best point estimate \emph{and}
a 90\% confidence interval (low, high).

\begin{enumerate}[resume, label=Q\arabic*.]

\item By the end of 2027, what do you think the highest reported score
      on GAIA will be? (Human baseline is $\approx 92\%$.)

\item By the end of 2027, what do you think the gap between the best
      LLM agent and humans on WebArena will be, in percentage points?
      (At April 2026: agent 71.6\%, human $\approx 78\%$, so gap
      $\approx 6.4$ points.)

\item Imagine a ``clean'' agentic-AI benchmark released in 2026 that
      has \emph{not} yet been contaminated. By the end of 2028, what
      fraction of the published leaderboard entries on that benchmark
      do you expect to show evidence of contamination?

\item Of the three ``grey zone'' preprints cited in this lecture
      (Sleeper Agents, In-Context Scheming, METR autonomous tasks),
      how many do you think will be published in a peer-reviewed venue
      by the end of 2027?

\end{enumerate}

\section*{Part III --- Scenarios}

\begin{enumerate}[resume, label=Q\arabic*.]

\item A vendor demonstrates an agent that reads inbound customer
      emails, drafts replies, and schedules follow-up meetings. The
      demo goes well. Which two of the four failure modes from this
      lecture are most relevant to the risk assessment for this
      deployment, and what is one concrete question you would ask the
      vendor about each?

\item Your CTO circulates a slide citing a SWE-bench score of 78\% for
      a product you are considering, and concludes ``this model can now
      fix most engineering tickets autonomously.'' Write a two-sentence
      reply that challenges the implicit claim using specific points
      from this lecture.

\item You are writing your procurement template. Pick three of the
      seven procurement questions from the deck. For each, write one
      concrete follow-up question you would use in an interview with
      the vendor.

\end{enumerate}

\vspace{2em}

\noindent\textbf{Scoring note.} Part~I is objective (one correct
answer). Part~II has no ``right'' answer; the useful output is your
own calibration over time. Part~III is discussion-oriented and is
best reviewed with a colleague or in a reading group.

\end{document}
